% !TeX root = ../tfg.tex
% !TeX encoding = utf8
%
%*******************************************************
% Capítulo 4. Formalización de la detección OOD y métricas
%*******************************************************

\chapter{Formalización de la detección OOD y métricas}\label{cap:ood-metricas}

\section{Distribución de los datos y cambios de distribución}\label{sec:cambios-distribucion}

% - Factorización P(X, Y) = P(Y | X) P(X).
% - Desplazamiento de covariables: cambia P(X), se mantiene P(Y | X).
% - Cambio semántico: cambia el espacio de etiquetas, P(Y) \neq P'(Y).
\begin{definicion}[Desplazamiento de covariables]\label{def:covariate-shift}
\pendiente{definición}
\end{definicion}

\begin{definicion}[Cambio semántico]\label{def:semantic-shift}
\pendiente{definición}
\end{definicion}

\section{Regla de decisión por umbral}\label{sec:regla-umbral}

% - Puntuación s: \mathcal{X} \to \R (más alto = más ID) y regla G_\tau(x) = ID si s(x) \geq \tau.
% - Convención de OpenOOD y de este trabajo: clase positiva = OOD.
\begin{definicion}[Detector OOD por umbral]\label{def:detector}
\pendiente{definición}
\end{definicion}

\section{Curva ROC y AUROC}\label{sec:auroc}

\begin{definicion}[Curva ROC]\label{def:roc}
% TPR(\tau) y FPR(\tau) como funciones del umbral; curva paramétrica (Fawcett, 2006).
\pendiente{definición}
\end{definicion}

\begin{teorema}[Interpretación probabilística del AUROC]\label{teo:auroc-prob}
% \AUROC = \Prob(s(X_{ID}) > s(X_{OOD})) (+ 1/2 \Prob(empate)) con X_{ID} y X_{OOD} independientes.
\pendiente{enunciado}
\end{teorema}
\begin{demostracion}
\pendiente{demostración}
\end{demostracion}

\begin{corolario}[AUROC empírico y estadístico U]\label{cor:auroc-u}
% \widehat{\AUROC} = U / (n_1 n_2), con U el estadístico de Mann-Whitney (\cref{sec:mann-whitney}).
\pendiente{enunciado}
\end{corolario}
\begin{demostracion}
\pendiente{demostración}
\end{demostracion}

\section{FPR@95}\label{sec:fpr95}

\begin{definicion}[FPR@TPR$x$]\label{def:fpr95}
\pendiente{definición}
\end{definicion}

\section{AUPR-In y AUPR-Out}\label{sec:aupr}

\begin{definicion}[Precisión, exhaustividad y AUPR]\label{def:aupr}
\pendiente{definición}
\end{definicion}

\section{CCR@FPR$x$}\label{sec:ccr}

\begin{definicion}[CCR@FPR$x$]\label{def:ccr}
\pendiente{definición}
\end{definicion}

\endinput
